\subsection{Элементарные преобразования матриц. Матрицы элементарных преобразований. Теоремы о матрицах элементарных преобразований}

\begin{definition}
Пусть $A$ — прямоугольная матрица размера $m \times n$. Элементарными преобразованиями строк называются действия следующего вида:
\begin{enumerate}
    \item Перестановка двух строк матрицы.
    \item Умножение строки на число $\lambda \neq 0$.
    \item Прибавление к одной строке другой строки этой же матрицы, умноженной на произвольное число $\beta$.
\end{enumerate}
Аналогично определяются элементарные преобразования столбцов.
\end{definition}

Данные действия являются обратимыми. Это означает, что если из матрицы $A$ путем элементарного преобразования получена матрица $A'$, то существует обратное преобразование того же типа, переводящее $A'$ в $A$:
\begin{itemize}
    \item Для перестановки строк $i$ и $j$ обратным действием является повторная перестановка тех же строк.
    \item Для умножения строки на $\lambda \neq 0$ обратным действием является умножение этой строки на $\frac{1}{\lambda}$.
    \item Для прибавления к $i$-й строке $j$-й строки, умноженной на $\beta$, обратным действием является вычитание из новой $i$-й строки $j$-й строки, умноженной на ту же $\beta$.
\end{itemize}

Для обозначения того, что матрицы получены друг из друга элементарными преобразованиями, используется знак эквивалентности: $A \sim \tilde{A}$.

\begin{definition}
Матрицами элементарных преобразований называются квадратные матрицы размера $m \times m$, полученные из единичной матрицы $E$ одним элементарным преобразованием:
\begin{enumerate}
    \item $B_{ij}$ — матрица перестановки: в единичной матрице на $i$-й строке единица сдвигается на $j$-е место, а на $j$-й строке единица сдвигается на $i$-е место.
    \[
    (B_{ij})_{sk} = \begin{cases} 1, & s=i, k=j \text{ или } s=j, k=i \\ 1, & s=k, s \neq i, j \\ 0, & \text{в остальных случаях} \end{cases}
    \]
    \item $B_{\lambda, i}$ — матрица умножения: в единичной матрице на $i$-м месте главной диагонали вместо единицы стоит число $\lambda \neq 0$.
    \[
    (B_{\lambda, i})_{sk} = \begin{cases} \lambda, & s=k=i \\ 1, & s=k \neq i \\ 0, & s \neq k \end{cases}
    \]
    \item $B_{i+\beta j}$ — матрица прибавления: берется единичная матрица, и в $i$-ю строку на $j$-е место добавляется число $\beta$. Таким образом, в $i$-й строке два ненулевых элемента: единица на диагонали и $\beta$ на $j$-м месте.
    \[
    (B_{i+\beta j})_{sk} = \begin{cases} 1, & s=k \\ \beta, & s=i, k=j \\ 0, & \text{в остальных случаях} \end{cases}
    \]
\end{enumerate}
\end{definition}

\begin{theorem}
Каждое элементарное преобразование строк матрицы $A$ сводится к умножению матрицы $A$ слева на соответствующую матрицу элементарных преобразований.
\end{theorem}

\begin{tcolorbox}[title=Доказательство, breakable]
Рассмотрим произведение $C = B A$. Элементы матрицы $C$ вычисляются по формуле:
\[ C_{sk} = \sum_{l=1}^m B_{sl} A_{lk} \]
Проверим теорему для каждого типа преобразований:

\textbf{1. Перестановка строк $i$ и $j$ (используем матрицу $B_{ij}$):}
\begin{itemize}
    \item Если $s \neq i$ и $s \neq j$, то в $s$-й строке матрицы $B_{ij}$ только один ненулевой элемент $B_{ss} = 1$. Тогда $C_{sk} = B_{ss} A_{sk} = A_{sk}$, то есть строки с номерами $s \neq i, j$ не меняются.
    \item Если $s = i$, то в $i$-й строке матрицы $B_{ij}$ ненулевым является только элемент $B_{ij} = 1$ (на месте $l=j$). Тогда $C_{ik} = B_{ij} A_{jk} = A_{jk}$. Таким образом, новая $i$-я строка совпадает со старой $j$-й строкой.
    \item Если $s = j$, аналогично: ненулевым в строке является $B_{ji} = 1$ (на месте $l=i$). Тогда $C_{jk} = B_{ji} A_{ik} = A_{ik}$. Новая $j$-я строка совпадает со старой $i$-й строкой.
\end{itemize}
Следовательно, умножение на $B_{ij}$ переставляет $i$-ю и $j$-ю строки.

\textbf{2. Умножение $i$-й строки на $\lambda$ (используем матрицу $B_{\lambda, i}$):}
\begin{itemize}
    \item Если $s \neq i$, то $B_{ss} = 1$, остальные элементы в строке нули. $C_{sk} = A_{sk}$.
    \item Если $s = i$, то единственным ненулевым элементом в строке является $B_{ii} = \lambda$. Тогда $C_{ik} = B_{ii} A_{ik} = \lambda A_{ik}$.
\end{itemize}
Вся $i$-я строка умножилась на $\lambda$.

\textbf{3. Прибавление к $i$-й строке $j$-й строки, умноженной на $\beta$ (используем матрицу $B_{i+\beta j}$):}
\begin{itemize}
    \item Если $s \neq i$, то $B_{ss} = 1$, остальные нули. $C_{sk} = A_{sk}$.
    \item Если $s = i$, то в $i$-й строке два ненулевых элемента: $B_{ii} = 1$ и $B_{ij} = \beta$. Сумма принимает вид:
    \[ C_{ik} = B_{ii} A_{ik} + B_{ij} A_{jk} = 1 \cdot A_{ik} + \beta \cdot A_{jk} \]
\end{itemize}
То есть к элементам $i$-й строки прибавились элементы $j$-й строки, умноженные на $\beta$.
Теорема доказана.
\end{tcolorbox}

\begin{tcolorbox}[title=Источники, colback=gray!10]
\begin{itemize}
  \item Лекция №\,1, тема «Элементарные преобразования строк и столбцов», временной отрезок 04:00--36:00
  \item PDF «Линейная алгебра\_compressed (1).pdf», раздел 1 «Матрица элементарных преобразований», стр. 1--3
  \item PDF «Вопросы», вопрос №\,1
\end{itemize}
\end{tcolorbox}
